\begin{abstract}
    \label{sec: abstract}
    Mean Flows enable distillation-free one-step generation by learning two-parameter average-velocity fields from noise to data distributions, but their training is notorious for non-decreasing loss and high gradient variance. We establish a theoretical framework that explains the phenomenon through the lens of \textit{variance reduction with a control variate}. We identify two distinct statistical roles for the conditional velocity $\vv_{\text{cond}}$ in training: as an unbiased regression target and as the Monte Carlo control variate for the tangent in the Jacobi-vector-product (JVP), but with a statistically suboptimal coefficient. The latter inflates per-step gradient variance and produces the non-decreasing loss. We derive the optimal tangent-mixing coefficient $\beta^{\ast}$ in closed form and show that a family of concurrent works correspond to different practical realizations of the same optimum. To validate the framework, we instantiate the $\beta\!=\!1$ corner via an EMA-derived proxy with a flow-matching anchor and run a controlled $\beta$-sweep ($\beta\!\in\!\{0,0.25,0.5,1\}$) at both toy and DiT-B/4 / ImageNet-$256$ scale. The sweep recovers the predicted bias-variance ordering at every matched-step checkpoint, with $3\!\times\!\sim\!5\!\times$ gradient-variance reduction and up to $54\%$ SW$_{1}$ improvement on 2-D benchmarks, and $\approx\!5\!\times$ reduction in gradient noise ratio at DiT scale. The same DiT-scale measurement also exposes a quantitative \emph{FID-vs-MSE landscape mismatch}: the FID-axis bias amplification across $\beta\!=\!0.5/\beta\!=\!1$ ($\sim\!1.4{:}1$) is sublinear in the gradient-MSE bias coefficient ($1{:}4$), so the FID-optimal $\beta$ shifts beyond the gradient-MSE interior minimizer to the unbiased corner $\beta\!=\!0$.
\end{abstract}
